\section{The supported relative simplicial construction and relational homotopy}
\label{sec:supported-dold-kan-independent}

The original relative construction is due to Alain Connes and
Caterina Consani,
\href{https://arxiv.org/abs/2004.08879v1}{\emph{Segal's Gamma rings
and universal arithmetic}, arXiv:2004.08879v1},
original \nolinkurl{Froblifts.tex}, lines 875--1231.
We use Proposition \texttt{relativeH} and the interval construction
\texttt{DKsimple}. The other inputs are the bounded functor
\texttt{simpconstruct1}, the elementary relation \texttt{Maydefn},
and Proposition \texttt{prophigher}. Their original amplitude construction and
its relational interpretation are antecedents. The full additive
carrier is the coefficient construction of The Clankers,
\emph{Split Support Geometry}, version 11,
\href{https://github.com/KokunoYumeto/zeta-function-research-reader/blob/a2851f9eb352e7e4376bc629de86fa4ed9c1c434/workbenches/splitzero-tandem/foundations/split-support-geometry-v11/split_support_geometry_arithmetic_curve_v11.tex#L1226}{Definition \texttt{def:lattice-split}}.

The source displays \texttt{pidegn}, lines 1133--1135, and
\texttt{pidegnk}, lines 1219--1221, omit one amplitude slot.
Degree $n$ has $n$ $A$ slots and one $B$ slot, as its earlier
definition and degree-three example require. We derive the faces
from the interval construction below. Homotopy means the source's
sphere and elementary relation, not homotopy of a geometric
realization.

\subsection{The original divisor and full additive carriers}

Fix
\begin{equation}\tag{SDKR1}
\begin{gathered}
D=\sum_p n_p[p]+a[\infty],\qquad
\Lambda_D=H^0(\operatorname{Spec}\mathbb Z,\mathcal O(D_{\rm fin}))
       =u_D\mathbb Z,\qquad u_D=\prod_p p^{-n_p}>0,\\
c=e^a>0,\qquad A=\mathbb R,\qquad B=\mathbb R/\Lambda_D,\qquad
\phi:A\to B,\quad \phi(r)=r+\Lambda_D.
\end{gathered}
\end{equation}
The equality for $\Lambda_D$ follows from all the original
conditions $v_p(x)+n_p\ge0$: $x/u_D$ has every finite valuation
nonnegative and is therefore an integer. No divisor coefficient,
archimedean bound or period is changed.

Let $L$ be the original nontrivial bounded distributive lattice.
For an abelian group $C$ put
\begin{equation}\tag{SDKR2}
\begin{gathered}
C_L=\{(0,\ell):\ell\in L\}\cup\{(x,1_L):x\in C\},\\
(x,\ell)+(y,m)=(x+y,\ell\vee m),\qquad
\tau_C=(0,0_L),\qquad e_C=(0,1_L).
\end{gathered}
\end{equation}
The pair $(0,1_L)$ occurs once. If a sum has nonzero amplitude,
a nonzero input forces its support to be top, so the operation
is closed. It is an associative commutative monoid law, with
identity $\tau_C$; empty sums mean $\tau_C$.
The quotient map lifts additively as
\begin{equation}\tag{SDKR3}
\phi_L:A_L\to B_L,\qquad (r,\ell)\mapsto(\phi(r),\ell).
\end{equation}
A nonzero $r\in\Lambda_D$ maps to $e_B$, not $\tau_B$.
The construction uses the additive groups; no
$\mathbb R$-module structure on $B$ is assumed.

\subsection{The relative functor and its exact simplicial coordinates}

For a finite pointed pair $(X,Y)$, $*\in Y\subseteq X$, define
\begin{equation}\tag{SDKR4}
\begin{split}
\mathcal K_D^L(X,Y)=\biggl\{\psi:\ &
Y\setminus\{*\}\to B_L,\quad X\setminus Y\to A_L:\\
&\sum_{x\in X\setminus Y}|\operatorname{amp}\psi(x)|\le c
\biggr\}.
\end{split}
\end{equation}
Its basepoint is the all-$\tau$ function. For a map of pointed
pairs $f:(X,Y)\to(X',Y')$ its nonbasepoint coordinates are
\begin{equation}\tag{SDKR5}
(\mathcal K_D^L(f)\psi)(z)=
\begin{cases}
\displaystyle\sum_{x\in X\setminus Y,\ f(x)=z}\psi(x),
                              &z\in X'\setminus Y',\\[4pt]
\displaystyle\sum_{y\in Y\setminus\{*\},\ f(y)=z}\psi(y)
+\displaystyle\sum_{x\in X\setminus Y,\ f(x)=z}\phi_L(\psi(x)),
                              &z\in Y'\setminus\{*\}.
\end{cases}
\end{equation}
In each fibre amplitudes add and supports join. A point outside
$Y'$ receives only points outside $Y$. The triangle inequality
therefore bounds the sum of its output absolute amplitudes by
the original outside sum; coordinates entering $Y'$ or the
basepoint are omitted from this bound. Thus the output lies in
SDKR4.
For a composite, partition each output fibre into the intermediate
fibres. Associativity of sums, additivity of $\phi_L$, and
associativity of finite joins give exactly the same amplitudes and
supports in the two calculations. This argument includes empty
fibres. Identity maps have singleton fibres. Hence SDKR5 defines
a pointed covariant functor, with no use of additive inverses.
Apply it to the source's interval pair
$\partial[n]^*=(\{0,\ldots,n+1\},\{0,n+1\})$, pointed at $0$,
smashed with $k_+$. In degree $n$ it gives
\begin{equation}\tag{SDKR6}
\begin{split}
\mathcal X_n(D;k)=\biggl\{(a_1,\ldots,a_n;b):\
&a_i\in A_L^k,\quad b\in B_L^k,\\
&\sum_{i=1}^n\sum_{j=1}^k
        |\operatorname{amp}(a_{i,j})|\le c\biggr\}.
\end{split}
\end{equation}
Thus $\mathcal X_0(D;k)=B_L^k$. For $k=0$ every degree is
a singleton. The final endpoint is the $B$ slot and the initial
endpoint is discarded. SDKR5 gives
\begin{equation}\tag{SDKR7}
\begin{aligned}
d_0(a_1,\ldots,a_n;b)&=(a_2,\ldots,a_n;b),\\
d_i(a_1,\ldots,a_n;b)
 &=(a_1,\ldots,a_{i-1},a_i+a_{i+1},a_{i+2},\ldots,a_n;b),
 &&1\le i<n,\\
d_n(a_1,\ldots,a_n;b)
 &=(a_1,\ldots,a_{n-1};b+\phi_L(a_n)).
\end{aligned}
\end{equation}
In particular $d_0(a;b)=b$ and $d_1(a;b)=b+\phi_L(a)$.
The degeneracies are
\begin{equation}\tag{SDKR8}
s_i(a_1,\ldots,a_n;b)
=(a_1,\ldots,a_i,\tau_A^k,a_{i+1},\ldots,a_n;b),
\qquad 0\le i\le n.
\end{equation}
Indeed a collapsed interior fibre adds adjacent $A$ entries;
a fibre entering the last endpoint applies $\phi_L$ and adds
to $B$; and a degeneracy inserts an empty fibre, hence $\tau_A$.
The simplicial identities follow from the interval maps and
the proved functoriality of SDKR5. This gives $n$ $A$ slots
and one $B$ slot in every degree, including all the support data.
For a pointed map $f:k_+\to l_+$ put
\begin{equation}\tag{SDKR9}
(f_*x)_h=\sum_{\substack{1\le j\le k\\f(j)=h}}x_j,\qquad
f_*:\mathcal X_n(D;k)\to\mathcal X_n(D;l),\quad
(a_i;b)\mapsto(f_*a_i;f_*b).
\end{equation}
The triangle inequality proves the bound after merging or
discarding coordinates. Additivity of $\phi_L$ proves commutation
with SDKR7--8. Fibre partition proves composition. These are the
raw pointed $\Gamma$ maps with the all-$\tau$ basepoint.
\subsection{Spheres and elementary homotopy at every supported basepoint}
For $s=(s_j)_j\in L^k$ let $z_s=((0,s_j))_j\in B_L^k$.
Its constant simplices are
\begin{equation}\tag{SDKR10}
\kappa_s^{(n)}
=(\underbrace{\tau_A^k,\ldots,\tau_A^k}_{n\ {\rm slots}};z_s).
\end{equation}
SDKR7--8 prove their basepoint identities. An $n$-sphere
for $n\ge1$ has every face equal to $\kappa_s^{(n-1)}$.
The complete set of one-spheres is
\begin{equation}\tag{SDKR11}
\begin{split}
\mathcal S_1(D;k,s)=\biggl\{(a;z_s):\
&a_j=(r_j,\ell_j)\in(\Lambda_D)_L,\quad \ell_j\le s_j,\\
&\sum_j|r_j|\le c\biggr\}.
\end{split}
\end{equation}
In fact $d_0(a;b)=z_s$ forces $b=z_s$. The other face equation
is exactly $\phi(r_j)=0$ and $\ell_j\vee s_j=s_j$ in every
coordinate. These are precisely the displayed conditions,
and conversely they imply both faces. The carrier condition
remains: $r_j\ne0$ forces $\ell_j=s_j=1_L$.
The source elementary relation on these spheres is
\begin{equation}\tag{SDKR12}
x\mathcal R_{1,s}y\ \Longleftrightarrow\
\exists w\in\mathcal X_2(D;k):
d_0w=\kappa_s^{(1)},\quad d_1w=x,\quad d_2w=y.
\end{equation}
Write $w=(u,v;b)$. The first equation forces
$v=\tau_A^k$, $b=z_s$, including their actual support labels.
Its other two faces are both $(u;z_s)$. Thus the elementary
relation is equality. Conversely, for $x=y=(a;z_s)$ the
witness $(a,\tau_A^k;z_s)$ has the required faces and bound.
In particular distinct zero-amplitude support labels are
distinct relational one-sphere classes.
For all higher degrees,
\begin{equation}\tag{SDKR13}
\mathcal S_n(D;k,s)=\{\kappa_s^{(n)}\}\qquad(n\ge2).
\end{equation}
Indeed $d_0x=\kappa_s^{(n-1)}$ forces
$a_2=\cdots=a_n=\tau_A^k$ and $b=z_s$.
The last face then forces $a_1=\tau_A^k$.
The constant simplex satisfies every face and has its reflexive
elementary witness in the next constant simplex. This proves
triviality of the source's higher sphere relations directly.
The extremes are therefore
\begin{equation}\tag{SDKR14}
\begin{aligned}
\mathcal S_1(D;k,0_L^k)&=\{(\tau_A^k;\tau_B^k)\},\\
\mathcal S_1(D;k,1_L^k)
 &=\{(a;e_B^k):a\in(\Lambda_D)_L^k,\
                          \sum_j|\operatorname{amp}a_j|\le c\}.
\end{aligned}
\end{equation}
The first is the loop object at the actual pointed zero.
The second, at the supported basepoint, retains the complete
bounded $\Gamma$-set of global sections in SDM3. The two
basepoints give different results and have not been identified.
\subsection{The exact finite-lattice count}
Suppose $L$ is finite, put $N=\lfloor c/u_D\rfloor$, and let
$I_s=\{j:s_j=1_L\}$. Then
\begin{equation}\tag{SDKR15}
\#\mathcal S_1(D;k,s)=
\sum_{\substack{T\subseteq I_s\\ |T|\le N}}
 2^{|T|}\binom N{|T|}\prod_{j\notin T}|\downarrow s_j|,
\qquad
\downarrow s_j=\{\ell\in L:\ell\le s_j\}.
\end{equation}
To prove this, choose the exact set $T$ of nonzero amplitudes.
It must lie in $I_s$, and those supports are forced to be top.
Every other coordinate has zero amplitude and independently
chooses a label in $\downarrow s_j$. The nonzero amplitudes
are $u_Dm_j$ for nonzero integers $m_j$. For $m=|T|$
there are $2^m$ sign choices. The positive magnitudes sum to
at most $N$. Their successive partial sums give an increasing
$m$-element subset of $\{1,\ldots,N\}$, and taking differences
is the inverse. Hence there are exactly $\binom Nm$ magnitude
choices. The empty set has one empty sign and magnitude choice.
Multiplying and summing proves the formula.
At full support it becomes
\begin{equation}\tag{SDKR16}
\#\mathcal S_1(D;k,1_L^k)
=\sum_{m=0}^{\min(k,N)}
 |L|^{k-m}2^m\binom km\binom Nm,
\qquad N=\left\lfloor\frac{e^a}{u_D}\right\rfloor .
\end{equation}
The original $a$ and $u_D$ remain in the expression.
If $e^a<u_D$, every amplitude is zero but there are $|L|^k$
distinct classes. At support $0_L^k$ there is one class.
The source's amplitude count omits these independent label
choices; its old symmetry is not asserted for SDKR16.
\subsection{Support transport and a genuine pointed Gamma construction}
For a pointed map $f:k_+\to l_+$, define
\begin{equation}\tag{SDKR17}
(f_*s)_h=\bigvee_{\substack{1\le j\le k\\ f(j)=h}}s_j,
\end{equation}
with bottom for an empty fibre. The raw map SDKR9 sends
$\kappa_s$ to $\kappa_{f_*s}$. In particular full support
maps to top in nonempty fibres and bottom in empty fibres.
For the insertion $1_+\to2_+$, its $B$ basepoint becomes
$(e_B,\tau_B)$, not $(e_B,e_B)$. All-top basepoints are
therefore not natural for the raw maps.
The exact support-indexed objects that repair the identity
law are
\begin{equation}\tag{SDKR18}
\mathcal X_n^{\ge s}(D;k)=
\{(a_i;b)\in\mathcal X_n(D;k):
             \operatorname{supp}(b_j)\ge s_j\text{ for all }j\}.
\end{equation}
They are simplicial subobjects: only the last face changes
the $B$ slot, by joining more labels, and degeneracies leave
it unchanged. Their basepoints are $\kappa_s$.
All the spheres and elementary witnesses above have $B=z_s$,
so their sphere objects and relations are unchanged.
For $s\le t$ there is a pointed simplicial transition
\begin{equation}\tag{SDKR19}
T_{s,t}:\mathcal X^{\ge s}(D;k)\to\mathcal X^{\ge t}(D;k),
\qquad (a_i;b)\mapsto(a_i;b+z_t).
\end{equation}
No $A$ amplitude changes. Commutation with the last face
is $b+z_t+\phi_L(a_n)=b+\phi_L(a_n)+z_t$;
the other faces and degeneracies commute immediately.
It sends $z_s$ to $z_t$. On its stated domain $T_{s,s}=I$,
since each $B$ label is already at least $s$.
For $s\le t\le u$ one has $T_{t,u}T_{s,t}=T_{s,u}$.
Adding $z_s$ would not be the identity on the unrestricted
object; SDKR18 is what proves this required identity.
Form the category whose objects are $(k_+,s)$, $s\in L^k$,
and whose morphisms $(k_+,s)\to(l_+,t)$ are pointed maps
$f$ with $f_*s\le t$. Monotonicity and
$(gf)_*=g_*f_*$ on support vectors prove that this is a
category. Its action is
\begin{equation}\tag{SDKR20}
\Theta_{f;s,t}(a_i;b)=(f_*a_i;f_*b+z_t).
\end{equation}
For a subsequent $(g;t,u)$, the composite $B$ slot is
\[
g_*(f_*b+z_t)+z_u
=(gf)_*b+z_{g_*t\vee u}
=(gf)_*b+z_u,
\]
and its $A$ slots are $(gf)_*a_i$. Thus
\begin{equation}\tag{SDKR21}
\Theta_{g;t,u}\Theta_{f;s,t}=\Theta_{gf;s,u},
\qquad \Theta_{\mathrm{id};s,s}=I.
\end{equation}
The maps are pointed and simplicial by their construction.
On one-spheres their exact action is
\begin{equation}\tag{SDKR22}
(a;z_s)\longmapsto(f_*a;z_t).
\end{equation}
Its supports are at most $f_*s\le t$ and its absolute-amplitude
sum does not increase. For fixed $k$, upward transition
is the inclusion of the same $A$ tuple into the larger
permitted support set.

At full support every pointed $f:k_+\to l_+$ is permitted
with target $1_L^l$. Thus SDKR20--21 yield a genuine pointed
simplicial $\Gamma$-object whose loop object is naturally
\begin{equation}\tag{SDKR23}
\Pi^{\rm rel}_1(\mathcal X^{\ge1_L^\bullet}(D))(k_+)
 \cong
 \{(r_j,\ell_j)_j\in(\Lambda_D)_L^k:
                                  \sum_j|r_j|\le e^a\}.
\end{equation}
The $\Gamma$ operation sums amplitudes and joins labels over
each fibre, putting $\tau_A$ in an empty fibre, by SDKR22.
Its pointed zero is the all-$\tau_A$ tuple. This proves
naturality of the full bounded module rather than assuming
naturality of the all-top basepoints in the unrestricted object.

\subsection{The complete degree-zero directed relation}

A vertex is exactly
\begin{equation}\tag{SDKR24}
(b,\nu)\in B^k\times L^k,\qquad
b_j\ne0\Longrightarrow\nu_j=1_L .
\end{equation}
The source relation is the existence of one simplex with
$d_0$ the first vertex and $d_1$ the second. It is
\begin{equation}\tag{SDKR25}
(b,\nu)\mathcal R_0(b',\mu)
\ \Longleftrightarrow\
\nu\le\mu,\qquad
\sum_jd_{\Lambda_D}(b_j,b'_j)\le c,
\end{equation}
where the exact original quotient distance is
\begin{equation}\tag{SDKR26}
d_{\Lambda_D}(r+\Lambda_D,r'+\Lambda_D)
 =\min_{m\in\mathbb Z}|r-r'+m u_D|.
\end{equation}
The minimum exists since the absolute values tend to
infinity as $|m|\to\infty$; it is independent of representatives
by reindexing.

For necessity, an edge $(a;(b,\nu))$, $a_j=(r_j,\ell_j)$,
ends at $(b+\phi(r),\nu\vee\ell)$. The support order follows.
Each $r_j$ represents the amplitude difference, so the
distance sum is at most $\sum_j|r_j|\le c$.
Conversely assume SDKR25 and choose minimum-absolute-value
representatives $r_j$ of $b'_j-b_j$. If $\mu_j<1_L$,
admissibility gives $b'_j=0$ and
$\nu_j\le\mu_j<1_L$ gives $b_j=0$; choose $r_j=0$.
All other target labels are top. Hence every
$(r_j,\mu_j)$ is a valid $A_L$ element, its total absolute
amplitude is at most $c$, and the edge with these $A$
coordinates ends at $(b',\nu\vee\mu)=(b',\mu)$.
This proves both directions.

In particular the singular relational object retaining
both data is
\begin{equation}\tag{SDKR27}
\begin{gathered}
\mathcal V_k=\{(b,\nu)\in B^k\times L^k:
                              b_j\ne0\Rightarrow\nu_j=1_L\},\\
\mathcal R_0=\{((b,\nu),(b',\mu))\in\mathcal V_k^2:
                     \nu\le\mu,\
                     d_{\Lambda_D}^{(1)}(b,b')\le c\},\\
d_{\Lambda_D}^{(1)}(b,b')=\sum_jd_{\Lambda_D}(b_j,b'_j).
\end{gathered}
\end{equation}
The amplitude projection maps onto the source distance
relation: every amplitude edge lifts with both supports
top. Its full-support section preserves the relation
in both directions. The support section
$\nu\mapsto(0,\nu)$ embeds the preorder $L^k$ with its
exact relation, since its amplitude distance is zero.
The joint projection is injective by SDKR24.

This relation is reflexive but need not be symmetric:
$(0,\nu)\to(0,\mu)$ exists precisely when $\nu\le\mu$.
For $k\ge1$ and nontrivial $L$ it is never the complete
relation, regardless of $c$, because an $e_B$ coordinate
cannot move to $\tau_B$. It can fail transitivity even
on full support: if $c<u_D/4$, the successive first-coordinate
amplitudes $0,c,2c$ modulo $\Lambda_D$, all other amplitudes
zero, give two permitted edges but first-to-last distance
$2c>c$. The one-step relation is retained rather than
replaced by its transitive closure.

\subsection{Naturality, the amplitude diameter and directed reachability}

Every pointed map satisfies
\begin{equation}\tag{SDKR28}
d_{\Lambda_D}^{(1)}(f_*b,f_*b')
   \le d_{\Lambda_D}^{(1)}(b,b').
\end{equation}
Indeed minimum difference representatives $r_j$ sum to
a representative in each output fibre. Its distance is
at most that absolute sum, which is at most the sum
of the $|r_j|$ in that fibre. Summing over output
nonbasepoint fibres proves the inequality. Thus raw
$\Gamma$ maps preserve SDKR25, using monotonicity of
support pushforward. For SDKR20 the vertex map is
$(b,\nu)\mapsto(f_*b,f_*\nu\vee t)$; joining the same
$t$ also preserves support order, proving its relational
naturality.

In the full-top subobject the support order is automatic.
Its amplitude relation is complete exactly when
\begin{equation}\tag{SDKR29}
c\ge ku_D/2
\quad\Longleftrightarrow\quad k\le2e^{\deg D},
\qquad \deg D=a-\log u_D .
\end{equation}
The circle $\mathbb R/u_D\mathbb Z$ has diameter exactly
$u_D/2$: representatives in $[-u_D/2,u_D/2]$ give the
upper bound and the class of $u_D/2$ attains it.
Taking that class in all $k$ coordinates gives exact
sum diameter $ku_D/2$. This proves the inclusive
endpoint in SDKR29. The second equivalence uses
$e^{\deg D}=e^a/u_D$, without changing $D$.

The separate directed reachability relation, allowing
finitely many source edges, is
\begin{equation}\tag{SDKR30}
(b,\nu)\mathcal R_0^*(b',\mu)
          \quad\Longleftrightarrow\quad \nu\le\mu.
\end{equation}
Necessity follows from support order along a path.
For sufficiency first use the zero-amplitude edge
$(b,\nu)\to(b,\mu)$. At fixed support $\mu$ all
coordinates below top have zero amplitudes at both
ends. Choose real minimum difference representatives
$r_j$, zero in those coordinates, and put
$R=\sum_j|r_j|$. Since $c>0$, choose $N\ge1$ with
$R/N\le c$. The vertices
$(b+\phi(i r/N),\mu)$ for $i=0,\ldots,N$ are valid.
Each successive edge has $A$ tuple $(r_j/N,\mu_j)_j$
of total absolute amplitude $R/N$. Its endpoint
after $N$ steps is exactly $(b',\mu)$.
This proves the finite path. Mutual reachability
is therefore equality of supports, with classes
indexed by all of $L^k$ and induced order its
coordinatewise order. This auxiliary calculation
does not change SDKR25 or assert a realization-homotopy
interpretation.

\subsection{Principal-divisor transport in the original coordinates}

For $q\in\mathbb Q^\times$, let $D'=D+(q)$ with
$(q)=\sum_pv_p(q)[p]-\log|q|[\infty]$. Then
\begin{equation}\tag{SDKR31}
\begin{gathered}
\Lambda_{D'}=q^{-1}\Lambda_D,\qquad
u_{D'}=u_D/|q|,\qquad c'=c/|q|,\\
(r,\ell)\mapsto(q^{-1}r,\ell),\qquad
(b,\ell)\mapsto(q^{-1}b,\ell).
\end{gathered}
\end{equation}
The quotient map on $B$ is well defined because its
lattice is carried exactly to $q^{-1}\Lambda_D$.
Both maps commute with $\phi_L$, preserve every join,
and scale every absolute amplitude by $|q|^{-1}$.
They therefore give inverse isomorphisms of all the
relative, simplicial and support-indexed constructions,
with inverse multiplication by $q$. SDKR26 shows that
distance scales by the same factor, proving preservation
of SDKR25 in both directions. Finally
$c'/u_{D'}=c/u_D$, so the counts are unchanged.
This proves invariance without first replacing $D$
by another representative.

\subsection{The controlling arithmetic-site lattice map}

Retain the original tropical coefficient semifield $K$
with zero and the extension of ASL18--22,
\begin{equation}\tag{SDKR32}
K_L=K\otimes_{\mathbb B}L,\qquad
\iota_L:L\to K_L,\quad\ell\mapsto1_K\otimes\ell,\qquad
J_L=\iota_L(L).
\end{equation}
The target support carrier here is exactly $J_L$,
not unrestricted $K_L$. In $J_L$, addition and product
are the joins and meets of the original labels:
the tensor relations give
$\iota_L(\ell\vee m)=\iota_L(\ell)+\iota_L(m)$ and
$\iota_L(\ell\wedge m)=\iota_L(\ell)\iota_L(m)$.
These elements are multiplicatively idempotent.

For completeness the labels are not identified by this
map. The nonzero-indicator map $K\to\mathbb B$, taking
zero to zero and each nonzero coefficient to one,
is a semiring map: in the positive tropical semifield,
a sum is zero exactly when both inputs are zero, and a
product is zero exactly when one is zero. Its tensor
extension
\[
\rho_L:K\otimes_{\mathbb B}L\to L,\qquad
\rho_L\!\left(\sum_i a_i\otimes\ell_i\right)
       =\bigvee_{a_i\ne0_K}\ell_i
\]
is well defined by the tensor universal property and
satisfies $\rho_L\iota_L=I_L$. Thus $\iota_L$ is
injective and identifies $L$ with the bounded
distributive sublattice $J_L$. This also agrees with
the joint-character reconstruction in ASL20--22:
a lattice character $\chi:L\to\mathbb B$ sends
$\iota_L(\ell)$ to $\chi(\ell)1_K$, and those
characters separate the original labels.

For either original additive group $C=A$ or $C=B$,
form its carrier with support restricted to $J_L$.
The controlling coefficient map and its inverse are
\begin{equation}\tag{SDKR33}
\begin{aligned}
\jmath_C:C_L&\xrightarrow{\ \sim\ }C_{J_L},
 &(x,\ell)&\longmapsto(x,1_K\otimes\ell),\\
\jmath_C^{-1}(x,\iota_L(\ell))&=(x,\ell).
\end{aligned}
\end{equation}
This is an isomorphism of additive monoids: amplitude
sums are unchanged and $\iota_L$ preserves joins and
bottom. It retains all coefficients, signs and constants.
Its square with the original quotient map commutes
exactly:
\[
\jmath_B\phi_L(r,\ell)
 =(\phi(r),\iota_L(\ell))
 =\phi_{J_L}\jmath_A(r,\ell).
\]
Applying it coordinatewise to SDKR4 therefore gives
a bijection
$\mathcal K_D^L(X,Y)\to\mathcal K_D^{J_L}(X,Y)$
with precisely the same amplitude bound. In SDKR5,
the $B$ terms commute by that square and the finite
support joins commute by the lattice homomorphism.
This proves naturality for every map of relative
pointed pairs, including every empty fibre. It
therefore commutes with every face, degeneracy
and raw $\Gamma$ map in SDKR6--9. Its basepoint
is transported by $s\mapsto\iota_L(s)$.
Because $\iota_L$ preserves and reflects order,
it maps SDKR18 bijectively to its corresponding
support-indexed subobject. Finally
$\iota_L(\ell\vee t)=\iota_L(\ell)\vee\iota_L(t)$
proves exact commutation with SDKR19--22.

The pair of amplitude and encoded support reconstructs
the complete relative simplex: in each $A$ or $B$
coordinate it belongs to exactly
\begin{equation}\tag{SDKR34}
\{(x,h)\in C\times J_L:
                    x\ne0\Rightarrow h=1_K\otimes1_L\},
\end{equation}
and the inverse is $(x,h)\mapsto(x,\rho_L(h))$.
On the relative object the only additional condition
is the unchanged sum of the outside absolute amplitudes.
Thus SDKR34 proves both injectivity and the exact joint
image, rather than asserting that amplitude and support
are unrelated presentations. It carries SDKR11, SDKR25
and their maps to identical reconstructed relations.

More generally, for any bounded lattice homomorphism
$\theta:L\to L'$ the coordinate map
$(x,\ell)\mapsto(x,\theta(\ell))$ gives a natural
map of the relative functors. It is well defined,
since $\theta(1_L)=1_{L'}$, and additive since
$\theta$ preserves joins and bottom. It commutes
with $\phi$, and therefore with SDKR5. Its support
basepoint is $\theta(s)$, and
$\theta(f_*s)=f_*\theta(s)$, including empty
fibres. It preserves the order condition for
SDKR20 and obeys
\begin{equation}\tag{SDKR35}
\theta_*\Theta_{f;s,t}
 =\Theta_{f;\theta(s),\theta(t)}\theta_*,
\qquad
(\mathrm{id}_K\otimes\theta)\iota_L
                   =\iota_{L'}\theta .
\end{equation}
The first equality follows by evaluating its $A$
coordinates and its $B$ coordinate:
both give the pushed amplitude and support
$\theta(f_*\nu)\vee\theta(t)$.
The second follows on each pure tensor
$1_K\otimes\ell$. These formulas prove the full
coefficient and support-transport squares for
the controlling lattice map. No choice of one
Boolean branch or loss of a zero label is
needed for SDKR33--34; a noninjective specialization
$\theta$ has exactly its stated image and does
not inherit that bijectivity claim.
